\documentclass[11pt]{article}

% ---------- Paquetes ----------
\usepackage[utf8]{inputenc}
\usepackage[T1]{fontenc}
\usepackage[spanish]{babel}
\usepackage[margin=2.2cm]{geometry}
\usepackage[scaled]{helvet}
\renewcommand{\familydefault}{\sfdefault}
\usepackage{tabularx}
\usepackage{array}
\usepackage[table]{xcolor}
\usepackage[most]{tcolorbox}
\usepackage{titlesec}
\usepackage{fancyhdr}
\usepackage{graphicx}
\usepackage{pdflscape}
\usepackage{graphicx}
\usepackage{attachfile}
\usepackage{amsmath}

% ---------- Colores ----------
\definecolor{principal}{HTML}{2E6B5E}
\definecolor{suave}{HTML}{E8F1EE}
\definecolor{acento}{HTML}{E0A458}
\definecolor{gris}{HTML}{5F6B6A}
\definecolor{colR}{HTML}{4F9D8A}
\definecolor{colA}{HTML}{2E6B5E}
\definecolor{colC}{HTML}{E0A458}
\definecolor{colI}{HTML}{C9D6D3}


% ---------- Títulos de sección ----------
\titleformat{\section}
    {\Large\bfseries\color{principal}}
    {}{0pt}
    {}[\vspace{-0.6em}\textcolor{acento}{\rule{\linewidth}{1.2pt}}]
\titlespacing*{\section}{0pt}{1.4em}{0.3em}

% ---------- Encabezado y pie ----------
\pagestyle{fancy}
\fancyhf{}
\fancyhead[L]{\small\color{gris}Envejecer mejor}
\fancyhead[R]{\small\color{gris}Planificación del proyecto}
\fancyfoot[C]{\small\color{gris}\thepage}
\renewcommand{\headrulewidth}{0.4pt}
\renewcommand{\headrule}{\color{suave}\hrule height 0.8pt}

\setlength{\parindent}{0pt}
\setlength{\parskip}{0.7em}

% ---------- Celdas de la matriz RACI ----------
\newcommand{\R}{\cellcolor{colR}\textcolor{white}{\textbf{R}}}
\newcommand{\A}{\cellcolor{colA}\textcolor{white}{\textbf{A}}}
\newcommand{\AR}{\cellcolor{colA}\textcolor{white}{\textbf{A/R}}}
\newcommand{\C}{\cellcolor{colC}\textbf{C}}
\newcommand{\I}{\cellcolor{colI}\textbf{I}}

% ---------- Encabezado de tabla ----------
\newcommand{\encab}[1]{\textcolor{white}{\textbf{#1}}}

\begin{document}

% ---------- Portada ----------
\begin{tcolorbox}[
    enhanced, colback=principal, colframe=principal,
    arc=4pt, boxrule=0pt, left=14pt, right=14pt, top=14pt, bottom=14pt
]
    {\color{white}\small PLANIFICACIÓN DEL PROYECTO}\\[0.4em]
    {\color{white}\Huge\bfseries Envejecer mejor}\\[0.6em]
    {\color{suave}\large Tareas, tiempos, Matriz RACI y presupuesto}
\end{tcolorbox}

\vspace{0.6em}

{\color{gris}\textbf{Integrantes:} Ulysses, Maximiliano, Fabián, Valentina y Lucas.}

% ---------- Objetivo ----------
\begin{tcolorbox}[
    enhanced, colback=suave, colframe=suave,
    borderline west={4pt}{0pt}{acento},
    arc=2pt, left=12pt, right=12pt, top=8pt, bottom=8pt
]
    {\color{principal}\bfseries Objetivo del trabajo}\\[0.3em]
    Informar y generar conciencia sobre cómo llegar a la tercera edad de forma
    saludable, mostrando qué hábitos ayudan a prevenir las enfermedades y
    afecciones más comunes de la vejez. Para eso haremos afiches y folletos
    impresos, publicaciones en redes sociales y una página web informativa.
\end{tcolorbox}

% ---------- Fechas clave ----------
\begin{tcolorbox}[
    enhanced, colback=white, colframe=suave, boxrule=1pt,
    arc=2pt, left=12pt, right=12pt, top=8pt, bottom=8pt
]
    {\color{principal}\bfseries Fechas clave}\\[0.3em]
    \renewcommand{\arraystretch}{1.3}
    \begin{tabular}{@{}ll}
        Inicio del trabajo & miércoles 30 de septiembre \\
        Presentación de avances & lunes 2 de noviembre \\
        Término de las tareas & martes 17 de noviembre \\
        Holgura (5 días) & 18 al 22 de noviembre \\
        Entrega del informe final y propuesta & lunes 23 de noviembre \\
    \end{tabular}\\[0.3em]
    {\small\color{gris}Además, cada semana se mostrará un avance del trabajo.}
\end{tcolorbox}

% ---------- Listado de tareas ----------
\clearpage
\section{Listado de tareas}

\begin{center}
    \renewcommand{\arraystretch}{1.7}
    \arrayrulecolor{white}
    \setlength{\arrayrulewidth}{1.5pt}
    \begin{tabularx}{\textwidth}{>{\centering\arraybackslash}p{0.8cm}
                                  >{\raggedright\arraybackslash}p{5.4cm}
                                  >{\raggedright\arraybackslash}X}
        \rowcolor{principal}
        \encab{N°} & \encab{Tarea} & \encab{Descripción} \\
        \rowcolor{suave}
        1 & Definir el rango etario &
            Investigar en qué edades conviene enfocar el mensaje y elegir el grupo objetivo. \\
        \rowcolor{white}
        2 & Investigar hábitos y enfermedades &
            Reunir información confiable sobre enfermedades de la vejez y los hábitos que ayudan a prevenirlas. \\
        \rowcolor{suave}
        3 & Definir mensajes y formato &
            Elegir las ideas principales que queremos transmitir y cómo se verán en cada material. \\
        \rowcolor{white}
        4 & Diseñar afiches, folletos y publicaciones &
            Crear los 10 afiches y folletos impresos y las publicaciones para redes sociales. \\
        \rowcolor{suave}
        5 & Desarrollar la página web &
            Programar una página informativa con HTML, CSS y JavaScript. \\
        \rowcolor{white}
        6 & Preparar la presentación de avances &
            Armar la presentación del 2 de noviembre. \\
        \rowcolor{suave}
        7 & Revisar contenidos y fuentes &
            Corregir errores y comprobar que la información esté bien respaldada. \\
        \rowcolor{white}
        8 & Difundir el material &
            Repartir los afiches y folletos, publicar en redes y compartir la página web. \\
        \rowcolor{suave}
        9 & Redactar el informe final y la propuesta &
            Escribir el informe con los resultados del proyecto y la propuesta final. \\
    \end{tabularx}
\end{center}

% ---------- Estimación de tiempo ----------
\clearpage
\section{Estimación de tiempo de cada tarea}

Cada integrante puede dedicar entre 3 y 5 horas a la semana, es decir, el
grupo cuenta con 15 a 25 horas semanales. Entre el 30 de septiembre y el 17 de
noviembre son 7 semanas, lo que da entre 105 y 175 horas disponibles.

\begin{center}
    \renewcommand{\arraystretch}{1.7}
    \arrayrulecolor{white}
    \setlength{\arrayrulewidth}{1.5pt}
    \begin{tabularx}{\textwidth}{>{\raggedright\arraybackslash}X
                                  >{\centering\arraybackslash}p{3.4cm}
                                  >{\centering\arraybackslash}p{1.4cm}
                                  >{\centering\arraybackslash}p{2.2cm}}
        \rowcolor{principal}
        \encab{Tarea} & \encab{Fechas} & \encab{Días} & \encab{Horas} \\
        \rowcolor{suave}
        1. Definir el rango etario & 30 sep -- 6 oct & 7 & 6 \\
        \rowcolor{white}
        2. Investigar hábitos y enfermedades & 5 oct -- 18 oct & 14 & 20 \\
        \rowcolor{suave}
        3. Definir mensajes y formato & 14 oct -- 21 oct & 8 & 8 \\
        \rowcolor{white}
        4. Diseñar afiches, folletos y publicaciones & 21 oct -- 1 nov & 12 & 15 \\
        \rowcolor{suave}
        5. Desarrollar la página web & 21 oct -- 10 nov & 21 & 30 \\
        \rowcolor{white}
        6. Preparar la presentación de avances & 26 oct -- 1 nov & 7 & 8 \\
        \rowcolor{suave}
        7. Revisar contenidos y fuentes & 3 nov -- 10 nov & 8 & 10 \\
        \rowcolor{white}
        8. Difundir el material & 4 nov -- 15 nov & 12 & 8 \\
        \rowcolor{suave}
        9. Redactar el informe final y la propuesta & 5 nov -- 17 nov & 13 & 20 \\
        \rowcolor{principal}
        \encab{Total} & & & \encab{125 horas} \\
    \end{tabularx}
\end{center}

{\small\color{gris}Las 125 horas quedan dentro de lo que el grupo puede dedicar. Algunas tareas se hacen en paralelo.}

% ---------- Matriz RACI ----------
\clearpage
\section{Matriz RACI}

Cada tarea tiene una sola persona que responde por el resultado (A), que en
algunos casos también es quien la ejecuta (A/R).

\begin{center}
    \renewcommand{\arraystretch}{1.8}
    \setlength{\tabcolsep}{4pt}
    \arrayrulecolor{white}
    \setlength{\arrayrulewidth}{1.5pt}
    \begin{tabularx}{\textwidth}{>{\raggedright\arraybackslash}X
                                  *{5}{>{\centering\arraybackslash}p{1.75cm}}}
        \rowcolor{principal}
        \encab{Tarea}
            & \textcolor{white}{\small\textbf{Ulysses}}
            & \textcolor{white}{\small\textbf{Maximiliano}}
            & \textcolor{white}{\small\textbf{Fabián}}
            & \textcolor{white}{\small\textbf{Valentina}}
            & \textcolor{white}{\small\textbf{Lucas}} \\
        \rowcolor{suave}
        1. Definir el rango etario                   & \AR & \R  & \C  & \C  & \I  \\
        \rowcolor{white}
        2. Investigar hábitos y enfermedades         & \C  & \AR & \R  & \C  & \I  \\
        \rowcolor{suave}
        3. Definir mensajes y formato                & \C  & \C  & \C  & \AR & \R  \\
        \rowcolor{white}
        4. Diseñar afiches, folletos y publicaciones & \I  & \C  & \R  & \R  & \AR \\
        \rowcolor{suave}
        5. Desarrollar la página web                 & \C  & \AR & \R  & \I  & \C  \\
        \rowcolor{white}
        6. Preparar la presentación de avances       & \AR & \R  & \C  & \R  & \C  \\
        \rowcolor{suave}
        7. Revisar contenidos y fuentes              & \R  & \R  & \AR & \C  & \C  \\
        \rowcolor{white}
        8. Difundir el material                      & \R  & \C  & \R  & \AR & \R  \\
        \rowcolor{suave}
        9. Redactar el informe final y la propuesta  & \R  & \R  & \C  & \C  & \AR \\
    \end{tabularx}
\end{center}

\begin{center}
    \setlength{\tabcolsep}{6pt}
    \renewcommand{\arraystretch}{1.5}
    \arrayrulecolor{white}
    \begin{tabular}{cl@{\hspace{1.2em}}cl}
        \R & \small Responsable: hace la tarea &
        \A & \small Aprobador: responde por el resultado \\[2pt]
        \C & \small Consultado: opina antes de decidir &
        \I & \small Informado: se entera del avance \\
    \end{tabular}
\end{center}

% ---------- Carta Gantt ----------
\section{Carta Gantt}

La Carta Gantt se encuentra en el archivo Excel adjunto (\texttt{carta\_gantt.xlsx}),
con las mismas tareas, fechas y responsables de este documento.
El archivo de la Carta Gantt en Excel se encuentra en el Anexo.

\begin{center}
    \includegraphics[width=\textwidth]{gantt.jpeg}
\end{center}

% ---------- Presupuesto ----------
\clearpage
\section{Presupuesto}

Valores referenciales en pesos chilenos. El presupuesto considera dos partes:
las horas hombre (HH) del equipo, valorizadas, y los gastos en materiales.

\textbf{Valor de la hora hombre.} Se usa como referencia el ingreso mínimo
mensual vigente desde el 1 de mayo de 2026, de \$553.553 para una jornada de
42 horas semanales (Ley N.º 21.830). El valor de la hora se calcula con la
fórmula de la Dirección del Trabajo:
\[
\text{Valor HH} = \frac{\$553.553}{30} \times \frac{28}{42} \approx \$12.301
\]

\subsection*{Horas hombre por tarea}

\begin{center}
    \renewcommand{\arraystretch}{1.6}
    \arrayrulecolor{white}
    \setlength{\arrayrulewidth}{1.5pt}
    \begin{tabularx}{\textwidth}{>{\raggedright\arraybackslash}X
                                  >{\centering\arraybackslash}p{1.6cm}
                                  >{\raggedleft\arraybackslash}p{2.4cm}
                                  >{\raggedleft\arraybackslash}p{2.6cm}}
        \rowcolor{principal}
        \encab{Tarea} & \encab{HH} & \encab{Valor HH} & \encab{Total} \\
        \rowcolor{suave}
        1. Definir el rango etario & 6 & \$12.301 & \$73.806 \\
        \rowcolor{white}
        2. Investigar hábitos y enfermedades & 20 & \$12.301 & \$246.020 \\
        \rowcolor{suave}
        3. Definir mensajes y formato & 8 & \$12.301 & \$98.408 \\
        \rowcolor{white}
        4. Diseñar afiches, folletos y publicaciones & 15 & \$12.301 & \$184.515 \\
        \rowcolor{suave}
        5. Desarrollar la página web & 30 & \$12.301 & \$369.030 \\
        \rowcolor{white}
        6. Preparar la presentación de avances & 8 & \$12.301 & \$98.408 \\
        \rowcolor{suave}
        7. Revisar contenidos y fuentes & 10 & \$12.301 & \$123.010 \\
        \rowcolor{white}
        8. Difundir el material & 8 & \$12.301 & \$98.408 \\
        \rowcolor{suave}
        9. Redactar el informe final y la propuesta & 20 & \$12.301 & \$246.020 \\
        \rowcolor{principal}
        \encab{Total horas hombre} & \encab{125} & & \encab{\$1.537.625} \\
    \end{tabularx}
\end{center}

\subsection*{Materiales y otros gastos}

\begin{center}
    \renewcommand{\arraystretch}{1.6}
    \arrayrulecolor{white}
    \setlength{\arrayrulewidth}{1.5pt}
    \begin{tabularx}{\textwidth}{>{\raggedright\arraybackslash}X
                                  >{\centering\arraybackslash}p{2cm}
                                  >{\raggedleft\arraybackslash}p{2.6cm}
                                  >{\raggedleft\arraybackslash}p{2.4cm}}
        \rowcolor{principal}
        \encab{Ítem} & \encab{Cantidad} & \encab{Valor unitario} & \encab{Total} \\
        \rowcolor{suave}
        Impresión a color de afiches y folletos & 10 & \$1.000 & \$10.000 \\
        \rowcolor{white}
        Página web (HTML, CSS y JavaScript, publicada en un hosting gratuito) & 1 & \$0 & \$0 \\
        \rowcolor{suave}
        Publicaciones en redes sociales & -- & \$0 & \$0 \\
        \rowcolor{white}
        Transporte & -- & \$0 & \$0 \\
        \rowcolor{principal}
        \encab{Total materiales} & & & \encab{\$10.000} \\
    \end{tabularx}
\end{center}

\subsection*{Resumen}

\begin{center}
    \renewcommand{\arraystretch}{1.6}
    \arrayrulecolor{white}
    \setlength{\arrayrulewidth}{1.5pt}
    \begin{tabularx}{0.75\textwidth}{>{\raggedright\arraybackslash}X
                                      >{\raggedleft\arraybackslash}p{3cm}}
        \rowcolor{principal}
        \encab{Concepto} & \encab{Monto} \\
        \rowcolor{suave}
        Horas hombre (125 HH) & \$1.537.625 \\
        \rowcolor{white}
        Materiales y otros gastos & \$10.000 \\
        \rowcolor{suave}
        Subtotal & \$1.547.625 \\
        \rowcolor{white}
        Imprevistos (10\,\% del subtotal) & \$154.763 \\
        \rowcolor{principal}
        \encab{Total del proyecto} & \encab{\$1.702.388} \\
    \end{tabularx}
\end{center}

{\small\color{gris}Las horas hombre son una valorización del trabajo del
equipo y no implican un pago. El gasto real en dinero es solo el de
materiales más sus imprevistos (\$10.000 + \$1.000 = \$11.000), por lo que
cada integrante aportaría aproximadamente \$2.200.}

% ---------- Anexo ----------
\clearpage
\section{Anexo: Carta Gantt}

Archivo original en Excel:
\textattachfile{carta_gantt.xlsx}{\textcolor{principal}{\textbf{carta\_gantt.xlsx}}}
(haz clic para abrirlo).
\end{document}
